\chapter{ĐÁNH GIÁ VÀ TRIỂN KHAI HỆ THỐNG}
\label{chap:eval}

Chương này trình bày việc đánh giá và triển khai hệ thống \IB{}: thiết lập đánh giá (Mục~\ref{sec:eval_setup}); đánh giá truy xuất (Mục~\ref{sec:eval_retrieval}); đánh giá chất lượng câu trả lời đầu cuối (Mục~\ref{sec:eval_answer}); thí nghiệm ablation theo từng thành phần (Mục~\ref{sec:ablation}); hiệu năng và chi phí (Mục~\ref{sec:perf}); triển khai hệ thống (Mục~\ref{sec:deploy}) và thảo luận (Mục~\ref{sec:discussion}). Các bảng số liệu trong chương được lấy từ tệp kết quả của đợt đánh giá ngày 24/11/2026 trên cùng kho tri thức và cùng bộ 200 câu hỏi: đợt này gồm 7 lượt đo mới và 4 lượt đo trước đó (hai cấu hình hybrid không xếp hạng lại và hai cấu hình dùng API Jina) được chấm lại theo cùng rubric; Precision@5 của các tệp truy xuất cũ được tính lại từ kết quả xếp hạng đã lưu.

% ==============================================================================
\section{Thiết lập đánh giá}
\label{sec:eval_setup}

\subsection{Bộ dữ liệu đánh giá}
\label{sec:golden_set}

Nhóm xây dựng một bộ câu hỏi đánh giá (golden set) gồm 200 câu, mỗi câu có đáp án tham chiếu, các URL nguồn đúng (gold) trên wiki và mức độ khó. Bộ câu được thiết kế để phản ánh câu hỏi thực tế của người chơi Việt Nam: hơn một nửa là tiếng Việt, có câu gõ sai, câu không dấu, câu viết tắt, câu nhiều lượt cần hỏi lại và câu ngoài phạm vi. Bảng~\ref{tab:golden_dist} trình bày phân bố theo độ khó và ngôn ngữ; Bảng~\ref{tab:golden_tags} trình bày các nhãn kiểu đầu vào.

\begin{table}[H]
    \centering
    \caption{Phân bố bộ câu hỏi đánh giá theo độ khó và ngôn ngữ}
    \label{tab:golden_dist}
    \small
    \begin{tabular}{|l|r|r|r|}
        \toprule
        \textbf{Độ khó} & \textbf{Tiếng Việt} & \textbf{Tiếng Anh} & \textbf{Tổng} \\
        \midrule
        Dễ & 33 & 27 & 60 \\ \hline
        Trung bình & 49 & 41 & 90 \\ \hline
        Khó & 28 & 22 & 50 \\ \hline
        \textbf{Tổng} & \textbf{110} & \textbf{90} & \textbf{200} \\
        \bottomrule
    \end{tabular}
\end{table}

\begin{table}[H]
    \centering
    \caption{Các nhãn kiểu đầu vào trong bộ câu hỏi đánh giá}
    \label{tab:golden_tags}
    \small
    \begin{tabularx}{\textwidth}{|>{\raggedright\arraybackslash}p{4.2cm}|>{\raggedright\arraybackslash}X|r|}
        \toprule
        \textbf{Nhãn} & \textbf{Ý nghĩa và ví dụ} & \textbf{Số câu} \\
        \midrule
        \code{typo} & Gõ sai tên riêng hoặc từ (Golden Sead, Gurrank, Death Night) & 37 \\ \hline
        \nolinkurl{vi_no_diacritics} & Tiếng Việt không dấu & 10 \\ \hline
        \code{abbreviation} & Viết tắt, teencode (k, j, z, bn, thg, coop) & 5 \\ \hline
        \code{telex} & Sót phím Telex (awn, giuwps, muoons) & 3 \\ \hline
        \code{casual\_case} & Viết thường, không dấu câu & 3 \\ \hline
        \code{clarification}, \code{multi\_turn} & Phải hỏi lại, có lượt bổ sung & 11 \\ \hline
        \code{out\_of\_scope} & Không liên quan đến Elden Ring & 6 \\
        \bottomrule
    \end{tabularx}
\end{table}

Một câu có thể mang nhiều nhãn; tổng cộng 49 câu có ít nhất một nhãn ``gõ không chuẩn'' (42 tiếng Việt, 7 tiếng Anh). Trong 200 câu có 183 câu dữ kiện (có đáp án và nguồn) và 17 câu hành vi (11 câu thiếu dữ kiện cần hỏi lại hoặc nêu các khả năng, 6 câu ngoài phạm vi). Phiếu review của mỗi câu dữ kiện lưu đoạn trích nguyên văn, mã đoạn, URL và mã băm của nguồn trong đúng bản dữ liệu dùng để đo. Bộ câu do thành viên phụ trách dữ liệu xây dựng và duyệt nguồn; lượt rà soát cuối trước khi đo được thực hiện với sự hỗ trợ của trợ lý AI theo ủy quyền của thành viên phụ trách đánh giá và được đánh dấu rõ trong dữ liệu là không phải duyệt trực tiếp bởi người.

\subsection{Rubric và chỉ số}
\label{sec:rubric}

Đơn vị chấm là một câu hỏi chính, gồm cả lượt bổ sung nếu có. Mỗi câu dữ kiện được chấm theo rubric ba mức:
\begin{itemize}
    \item Độ đúng (0/1/2): 2 điểm khi đủ ý bắt buộc, đúng chủ thể, điều kiện và ngôn ngữ; 1 điểm khi có ý đúng nhưng thiếu ý quan trọng hoặc thêm chi tiết sai không đảo ngược kết luận; 0 điểm khi sai, lỗi hoặc không trả lời.
    \item Độ bám nguồn (0/1/2): 2 điểm khi các khẳng định quan trọng được đoạn trích hỗ trợ trực tiếp; 1 điểm khi nguồn chỉ xác nhận một phần; 0 điểm khi khẳng định chính không có căn cứ.
    \item Hành vi (đạt/không đạt) cho câu hành vi: câu ngoài phạm vi đạt khi được nhận diện và từ chối phù hợp; câu thiếu dữ kiện đạt khi hệ thống hỏi đúng thông tin cần để phân biệt, hoặc trả lời có điều kiện với ít nhất hai khả năng. Chỉ đưa ra một khả năng là không đạt.
\end{itemize}
Từ điểm chấm, các chỉ số chính được tính như sau: \textbf{độ đúng hoàn toàn} (accuracy) là tỉ lệ câu dữ kiện đạt 2 điểm độ đúng trên toàn bộ 183 câu dữ kiện, kể cả câu lỗi; \textbf{faithfulness 2} là tỉ lệ câu dữ kiện đạt 2 điểm bám nguồn; \textbf{trích đúng nguồn gold} là tỉ lệ câu dữ kiện có trích dẫn ít nhất một trang gold; \textbf{behavior pass} là tỉ lệ câu hành vi đạt. Chỉ số truy xuất (Recall@5, MRR@5, Hit@5, Precision@5) được tính theo định nghĩa ở Mục~\ref{sec:ir_metrics} trên năm trang khác nhau đứng đầu, với câu hỏi nguyên văn. Chỉ số vận hành gồm độ trễ tới câu trả lời đầu tiên (p50, p95), số câu vượt 45 giây, số lần gọi công cụ và số token mỗi câu.

Người chấm. Câu trả lời được chấm bởi một mô hình ngôn ngữ (Claude, chạy qua Claude Code) theo rubric và quy ước chấm viết sẵn, không dùng chính DeepSeek để tránh thiên lệch tự chấm~\cite{zheng2023judge}. Để kiểm tra độ tin cậy, nhóm so sánh với một bộ điểm độc lập do mô hình \code{deepseek-v4-pro} chấm trên cấu hình hybrid không xếp hạng lại, web tự động: hai người chấm khớp điểm độ đúng ở 181/183 câu và cho cùng độ đúng hoàn toàn 95,6\%.

\subsection{Cấu hình đo}
\label{sec:eval_config}

\begin{itemize}
    \item Mô hình: \code{deepseek-flash} cho cả Agent và bước sinh; Agent tắt chế độ suy nghĩ; bước sinh được phép suy nghĩ khi Agent yêu cầu.
    \item Kho tri thức: collection 34.138 điểm mô tả ở Bảng~\ref{tab:corpus_stats}.
    \item Cách chạy: mỗi cấu hình chạy đủ 200 câu qua API thật (SSE), tuần tự từng câu, lưu hội thoại trong bộ nhớ, nâng giới hạn tần suất để không bị chặn. Không chạy song song vì các cấu hình có bộ xếp hạng lại cục bộ sẽ tranh CPU và vượt thời hạn 3 giây, làm sai lệch số đo.
    \item Các trục thay đổi: kiểu truy xuất (dense E5, BM25, hybrid RRF); bộ xếp hạng lại (không có, MiniLM-L12 cục bộ, API Jina v2 làm số tham khảo); chế độ web (tắt, tự động, luôn tìm). Tổng cộng 11 cấu hình, mỗi cấu hình chỉ đổi một trục so với cấu hình lân cận.
    \item Cấu hình mặc định của hệ thống: hybrid RRF + MiniLM-L12 + web tự động.
\end{itemize}

% ==============================================================================
\section{Đánh giá truy xuất}
\label{sec:eval_retrieval}

Truy xuất được đo độc lập với bước sinh trên 183 câu dữ kiện, dùng câu hỏi nguyên văn (không qua Agent viết lại), để đánh giá riêng chất lượng chỉ mục và bộ xếp hạng. Kết quả ở Bảng~\ref{tab:retrieval} và Hình~\ref{fig:retrieval_chart}.

\begin{table}[H]
    \centering
    \caption{Kết quả truy xuất trên 183 câu dữ kiện (câu hỏi nguyên văn)}
    \label{tab:retrieval}
    \small
    \begin{tabularx}{\textwidth}{|l|>{\raggedright\arraybackslash}X|r|r|r|r|r|}
        \toprule
        \textbf{Truy xuất} & \textbf{Xếp hạng lại} & \textbf{P@5} & \textbf{Recall@5} & \textbf{Hit@5} & \textbf{MRR@5} & \textbf{Thời gian} \\
        \midrule
        Dense E5 & -- & 0,189 & 81,1\% & 85,2\% & 0,714 & 38 ms \\ \hline
        Dense E5 & MiniLM-L12 & 0,201 & 85,5\% & 90,2\% & 0,733 & 0,74 s \\ \hline
        Dense E5 & API Jina v2 (tham khảo) & 0,208 & 87,7\% & 91,8\% & 0,769 & 2,0 s \\ \hline
        BM25 & -- & 0,183 & 78,6\% & 82,5\% & 0,584 & 37 ms \\ \hline
        BM25 & MiniLM-L12 & 0,192 & 82,5\% & 88,0\% & 0,718 & 0,75 s \\ \hline
        BM25 & API Jina v2 (tham khảo) & 0,210 & 88,3\% & 91,8\% & 0,763 & 2,3 s \\ \hline
        Hybrid RRF & -- & 0,195 & 83,1\% & 87,4\% & 0,706 & 67 ms \\ \hline
        \textbf{Hybrid RRF} & \textbf{MiniLM-L12 (mặc định)} & \textbf{0,203} & \textbf{86,6\%} & \textbf{90,7\%} & \textbf{0,739} & 1,0 s \\ \hline
        Hybrid RRF & API Jina v2 (tham khảo) & 0,215 & 90,7\% & 94,0\% & 0,789 & 2,0 s \\
        \bottomrule
    \end{tabularx}
\end{table}
\tablenote{Thời gian là trung bình mỗi truy vấn. Thời gian của các lượt API Jina gồm cả thời gian chờ hạn mức gọi API nên không phản ánh độ trễ thực của một lần gọi.}

\begin{figure}[H]
\centering
\begin{tikzpicture}
\begin{axis}[
    ybar, bar width=9pt, width=0.95\textwidth, height=6.2cm,
    ymin=50, ymax=100, ylabel={Recall@5 (\%)},
    symbolic x coords={Dense E5, BM25, Hybrid RRF},
    xtick=data, enlarge x limits=0.25,
    legend style={at={(0.5,-0.18)}, anchor=north, legend columns=3, font=\small},
    nodes near coords, nodes near coords style={font=\scriptsize, /pgf/number format/use comma},
    every node near coord/.append style={rotate=90, anchor=west},
    tick label style={font=\small}, label style={font=\small},
    /pgf/number format/use comma
]
\addplot[fill=AccentGray!50] coordinates {(Dense E5,81.1) (BM25,78.6) (Hybrid RRF,83.1)};
\addplot[fill=AccentBlue!60] coordinates {(Dense E5,85.5) (BM25,82.5) (Hybrid RRF,86.6)};
\addplot[fill=AccentLightBlue] coordinates {(Dense E5,87.7) (BM25,88.3) (Hybrid RRF,90.7)};
\legend{Không xếp hạng lại, MiniLM-L12 cục bộ, API Jina v2 (tham khảo)}
\end{axis}
\end{tikzpicture}
\caption{Recall@5 theo kiểu truy xuất và bộ xếp hạng lại}
\label{fig:retrieval_chart}
\end{figure}

Một số nhận xét:
\begin{itemize}
    \item Truy xuất lai tốt hơn từng phương pháp đơn lẻ (RQ1). Không xếp hạng lại, hybrid RRF đạt Recall@5 83,1\%, cao hơn dense E5 (81,1\%) và BM25 (78,6\%). Về MRR@5, hybrid (0,706) xấp xỉ dense (0,714) và cao hơn hẳn BM25 (0,584). Khi có MiniLM-L12, hybrid vẫn là tổ hợp tốt nhất trong các cấu hình không gọi API. Mọi cấu hình hybrid đều vượt yêu cầu Recall@5 $\geq$ 0,80.
    \item Bộ xếp hạng lại giúp mọi kiểu truy xuất. MiniLM-L12 tăng Recall@5 của hybrid từ 83,1\% lên 86,6\% và MRR@5 từ 0,706 lên 0,739, tức khoảng một nửa mức tăng của API Jina (90,7\% và 0,789) mà không tốn phí API. Mức tăng lớn nhất ở BM25: MRR@5 từ 0,584 lên 0,718, vì BM25 xếp hạng kém nhưng vẫn kéo được trang đúng vào nhóm 30 ứng viên đầu, để cross-encoder đẩy lên.
    \item Chi phí thời gian: truy xuất lai không xếp hạng lại chỉ mất khoảng 67 ms mỗi truy vấn; MiniLM-L12 trên CPU thêm khoảng 0,9 giây.
    \item Precision@5 ở mức 0,18--0,22 vì phần lớn câu chỉ có một đến hai trang gold; chỉ số này dùng để so sánh tương đối giữa các cấu hình.
\end{itemize}

% ==============================================================================
\section{Đánh giá chất lượng câu trả lời đầu cuối}
\label{sec:eval_answer}

\subsection{Kết quả của cấu hình mặc định}
\label{sec:main_results}

Với cấu hình mặc định (hybrid RRF + MiniLM-L12 + web tự động), cả 200/200 câu đều trả về thành công, không câu nào vượt hạn 45 giây. Trong 183 câu dữ kiện, 179 câu đạt độ đúng 2 điểm, 2 câu đạt 1 điểm và 2 câu 0 điểm. Bảng~\ref{tab:main_results} tóm tắt các chỉ số chính.

\begin{table}[H]
    \centering
    \caption{Kết quả đầu cuối của cấu hình mặc định trên bộ 200 câu}
    \label{tab:main_results}
    \small
    \begin{tabularx}{\textwidth}{|>{\raggedright\arraybackslash}X|r|>{\raggedright\arraybackslash}p{4.0cm}|}
        \toprule
        \textbf{Chỉ số} & \textbf{Kết quả} & \textbf{Yêu cầu đặt ra} \\
        \midrule
        Câu trả về thành công & 200/200 (100\%) & \\ \hline
        Độ đúng hoàn toàn (183 câu dữ kiện) & 97,8\% & $\geq$ 70\% \\ \hline
        Điểm độ đúng trung bình (thang 0--2) & 1,967 & \\ \hline
        Faithfulness 2 & 96,2\% & $\geq$ 85\% \\ \hline
        Trích đúng ít nhất một trang gold & 95,1\% & \\ \hline
        Behavior pass (17 câu hành vi) & 70,6\% (12/17) & \\ \hline
        Độ trễ tới câu trả lời đầu tiên p50 / p95 & 6,6 s / 15,1 s & p95 $\leq$ 10 s (mục tiêu) \\ \hline
        Số câu vượt 45 giây & 0 & 0 \\ \hline
        Số lần \code{search\_corpus} / \code{search\_web} mỗi câu & 2,23 / 0,01 & \\ \hline
        Token mỗi câu (Agent + bước sinh) & khoảng 22.100 & \\ \hline
        Câu phải kết thúc bằng lượt \code{finish} bắt buộc & 1 & \\
        \bottomrule
    \end{tabularx}
\end{table}

Hệ thống vượt các yêu cầu về độ đúng ($\geq$ 70\%) và độ bám nguồn ($\geq$ 85\%) với khoảng cách lớn. Yêu cầu độ trễ p95 $\leq$ 10 giây chưa đạt (p95 15,1 giây), dù không câu nào vượt hạn 45 giây của Discord bot.

\subsection{Kết quả theo nhóm câu hỏi}
\label{sec:group_results}

Bảng~\ref{tab:group_results} chia kết quả của cấu hình mặc định theo độ khó, ngôn ngữ và kiểu đầu vào (chỉ tính câu dữ kiện).

\begin{table}[H]
    \centering
    \caption{Kết quả của cấu hình mặc định theo nhóm câu dữ kiện}
    \label{tab:group_results}
    \small
    \begin{tabular}{|l|r|r|r|r|}
        \toprule
        \textbf{Nhóm} & \textbf{Số câu} & \textbf{Độ đúng hoàn toàn} & \textbf{Faithfulness 2} & \textbf{p50 / p95 (s)} \\
        \midrule
        Dễ & 55 & 96,4\% & 96,4\% & 5,4 / 10,4 \\ \hline
        Trung bình & 83 & 98,8\% & 96,4\% & 6,7 / 10,7 \\ \hline
        Khó & 45 & 97,8\% & 95,6\% & 10,5 / 19,4 \\ \hline
        Tiếng Việt & 97 & 96,9\% & 94,8\% & 7,3 / 15,8 \\ \hline
        Tiếng Anh & 86 & 98,8\% & 97,7\% & 6,2 / 14,5 \\ \hline
        Gõ không chuẩn & 49 & 98,0\% & 95,9\% & -- \\ \hline
        Gõ chuẩn & 134 & 97,8\% & 96,3\% & -- \\
        \bottomrule
    \end{tabular}
\end{table}
\tablenote{Độ trễ theo độ khó lấy từ tệp tổng hợp, tính trên tất cả câu của nhóm (gồm cả câu hành vi); độ trễ theo ngôn ngữ do nhóm tính từ kết quả từng câu dữ kiện, phân vị theo phương pháp nearest-rank.}

Kết quả cho thấy ba điểm. Thứ nhất, \textbf{câu gõ không chuẩn không còn là điểm yếu}: độ đúng của nhóm này (98,0\%) tương đương nhóm gõ chuẩn, nhờ Agent tự viết lại truy vấn tiếng Anh và kiểm tên bằng \code{check\_name}. Thứ hai, \textbf{câu tiếng Việt chỉ thấp hơn câu tiếng Anh khoảng hai điểm phần trăm} dù kho tri thức hoàn toàn tiếng Anh, cho thấy cách để Agent dịch truy vấn hoạt động tốt. Thứ ba, \textbf{câu khó tốn nhiều tài nguyên hơn}: trung bình 2,92 lần tìm kiếm và khoảng 29.400 token mỗi câu (câu dễ: 1,88 lần và khoảng 18.600 token), p95 độ trễ gần gấp đôi.

\subsection{Phân tích lỗi}
\label{sec:error_analysis}

Lỗi chủ yếu của hệ thống nằm ở hành vi với câu mơ hồ: 5 trong 17 câu hành vi không đạt, đều là câu thiếu dữ kiện. Các câu này lặp lại gần như giống nhau ở cả 11 cấu hình, nên không phụ thuộc vào truy xuất, xếp hạng lại hay tìm web. Bảng~\ref{tab:errors} liệt kê các câu fail thường xuyên nhất.

\begin{table}[H]
    \centering
    \caption{Các câu thiếu dữ kiện hệ thống thường xử lý chưa đạt}
    \label{tab:errors}
    \small
    \begin{tabularx}{\textwidth}{|>{\raggedright\arraybackslash}X|c|>{\raggedright\arraybackslash}X|}
        \toprule
        \textbf{Câu hỏi (rút gọn)} & \textbf{Số cấu hình fail} & \textbf{Hành vi của hệ thống} \\
        \midrule
        ``Vừa có remembrance đầu tiên, bước tiếp của quest NPC đó?'' & 11/11 & Mặc định là Remembrance of the Grafted và hướng dẫn tới gặp Enia \\ \hline
        ``Talisman rồng bản cộng đó tăng bao nhiêu kháng?'' & 11/11 & Tự chọn dòng Dragoncrest, không hỏi là dòng talisman rồng nào \\ \hline
        ``Hiệp sĩ ở Roundtable Hold, quest dẫn đi đâu?'' & 10/11 & Mặc định là Gideon Ofnir \\ \hline
        ``Rìu của kỵ sĩ chết và viên crystal tear xanh?'' & 10/11 & Liệt kê đủ hai cây rìu nhưng tự chọn một loại crystal tear \\ \hline
        ``Sau tiếng báo Great Rune vỡ trong DLC, gặp lại hiệp sĩ ở chỗ thánh giá?'' & 9/11 & Mặc định là Ansbach ở Main Gate Cross \\
        \bottomrule
    \end{tabularx}
\end{table}

Kiểu lỗi chung là câu hỏi có một tham chiếu mơ hồ (``hiệp sĩ đó'', ``bản cộng đó''); truy xuất tìm được một ứng viên hợp lý và Agent trả lời luôn kèm trích dẫn, thay vì hỏi lại hoặc liệt kê các khả năng. Đáng chú ý là các câu trả lời này không bịa: nội dung đúng với cách hiểu mà Agent chọn và có nguồn, nên độ đúng và độ bám nguồn của các câu dữ kiện không bị ảnh hưởng. Đây là vấn đề của chính sách hỏi lại trong prompt Agent, cụ thể là prompt đang ưu tiên ``hữu ích trước'' và chỉ hỏi lại khi câu trả lời thay đổi đáng kể. Ngược lại, toàn bộ 6 câu ngoài phạm vi được từ chối đúng.

Ngoài ra, ở hai cấu hình web ``luôn tìm'', một câu hỏi về việc dùng phép Carian Slicer với lớp Samurai trả lỗi HTTP 502 ở cả hai lần chạy, trong khi trả lời đúng ở chế độ tắt và tự động. Log cho thấy bước sinh đưa ra chỉ số bằng chứng vượt quá số đoạn ngữ cảnh, bị bộ kiểm tra từ chối ở cả lần sinh đầu và lần sinh lại. Bộ kiểm tra đã ngăn một câu trả lời trích dẫn sai đến tay người dùng, nhưng hệ quả là người dùng nhận thông báo lỗi; nguyên nhân gốc (nhiều khả năng do ngữ cảnh có thêm đoạn web khiến mô hình đánh số sai) chưa được xác định chắc chắn vì log không ghi nội dung mô hình.

% ==============================================================================
\section{Thí nghiệm ablation}
\label{sec:ablation}

Bảng~\ref{tab:ablation} trình bày kết quả đầu cuối của 11 cấu hình. Mỗi câu dữ kiện tương ứng khoảng 0,55 điểm phần trăm độ đúng, nên cấu hình tốt nhất và kém nhất chỉ chênh nhau 5 câu trên 183; phần lớn khác biệt giữa các cấu hình nằm trong phạm vi dao động ngẫu nhiên và cần được diễn giải thận trọng.

\begin{table}[H]
    \centering
    \caption{Kết quả đầu cuối của 11 cấu hình trên bộ 200 câu}
    \label{tab:ablation}
    \small
    \begin{tabularx}{\textwidth}{|>{\raggedright\arraybackslash}X|r|r|r|r|r|}
        \toprule
        \textbf{Cấu hình} & \textbf{Độ đúng} & \textbf{Faith. 2} & \textbf{Hành vi} & \textbf{Gold} & \textbf{p50 / p95 (s)} \\
        \midrule
        Dense E5, không xếp hạng lại, web tắt & 95,6\% & 96,7\% & 70,6\% & 95,1\% & 4,5 / 10,7 \\ \hline
        Dense E5 + L12, web tự động & 97,3\% & 97,8\% & 76,5\% & 95,6\% & 7,3 / 18,7 \\ \hline
        BM25 + L12, web tự động & 96,2\% & 98,4\% & 70,6\% & 94,5\% & 7,8 / 19,6 \\ \hline
        Hybrid RRF, web tắt & 95,1\% & 97,8\% & 70,6\% & 95,1\% & 5,0 / 12,8 \\ \hline
        Hybrid RRF, web tự động & 95,6\% & 97,8\% & 70,6\% & 92,9\% & 5,2 / 10,9 \\ \hline
        Hybrid RRF, web luôn tìm & 96,7\% & 99,5\% & 70,6\% & 94,5\% & 11,7 / 18,7 \\ \hline
        Hybrid + L12, web tắt & 96,7\% & 98,4\% & 70,6\% & 96,7\% & 6,6 / 15,0 \\ \hline
        \textbf{Hybrid + L12, web tự động (mặc định)} & \textbf{97,8\%} & 96,2\% & 70,6\% & 95,1\% & 6,6 / 15,1 \\ \hline
        Hybrid + L12, web luôn tìm & 96,7\% & 98,4\% & 70,6\% & 95,1\% & 14,6 / 25,7 \\ \hline
        Hybrid + API Jina, web tắt (tham khảo) & 96,2\% & 98,4\% & 76,5\% & 95,1\% & 5,9 / 14,0 \\ \hline
        Hybrid + API Jina, web tự động (tham khảo) & 96,7\% & 97,3\% & 70,6\% & 95,6\% & 6,0 / 13,1 \\
        \bottomrule
    \end{tabularx}
\end{table}
\tablenote{L12: cross-encoder MiniLM-L12 chạy cục bộ. Độ trễ là thời gian tới câu trả lời đầu tiên. Mọi cấu hình trả về 200/200 câu, trừ hai cấu hình web luôn tìm có 1 câu lỗi HTTP 502.}

\begin{figure}[H]
\centering
\begin{tikzpicture}
\begin{axis}[
    width=0.92\textwidth, height=7cm,
    xlabel={Độ trễ p50 tới câu trả lời đầu tiên (s)}, ylabel={Độ đúng hoàn toàn (\%)},
    xmin=3.5, xmax=16, ymin=94.5, ymax=98.5,
    grid=both, grid style={AccentGray!30},
    tick label style={font=\small}, label style={font=\small},
    /pgf/number format/use comma,
    nodes near coords, point meta=explicit symbolic,
    nodes near coords style={font=\scriptsize, anchor=south west, inner sep=1pt}
]
\addplot[only marks, mark=*, mark size=2.4pt, AccentBlue] coordinates {
  (4.5,95.6) [D-off]
  (7.3,97.3) [D+L12-auto]
  (7.8,96.2) [B+L12-auto]
  (5.0,95.1) [H-off]
  (5.2,95.6) [H-auto]
  (11.7,96.7) [H-always]
  (6.6,96.7) [H+L12-off]
  (14.6,96.7) [H+L12-always]
};
\addplot[only marks, mark=square*, mark size=3pt, red!70!black] coordinates {(6.6,97.8) [Mặc định: H+L12-auto]};
\end{axis}
\end{tikzpicture}
\caption{Đánh đổi giữa độ đúng và độ trễ của các cấu hình không gọi API xếp hạng lại}
\fignote{D: dense E5; B: BM25; H: hybrid RRF; off/auto/always: chế độ web tắt/tự động/luôn tìm.}
\label{fig:ablation_scatter}
\end{figure}

\subsection{Ảnh hưởng của bộ xếp hạng lại}
\label{sec:ablation_rerank}

Ở trọn luồng, MiniLM-L12 tăng độ đúng của hybrid khoảng 1--2 điểm phần trăm ở cả web tắt (95,1\% lên 96,7\%) và web tự động (95,6\% lên 97,8\%), đổi lại p50 tăng khoảng 1,5 giây. Tác dụng rõ nhất ở câu khó: mọi cấu hình không xếp hạng lại đều đúng 91,1\% câu khó, còn các cấu hình có xếp hạng lại ở web tắt và tự động đạt 93,3--97,8\%. Kết hợp với mức tăng Recall@5 và MRR@5 ở Mục~\ref{sec:eval_retrieval}, nhóm kết luận bộ xếp hạng lại cục bộ đáng giá với độ trễ tăng thêm (RQ2) và chọn làm mặc định. Bộ xếp hạng lại hoạt động thành công ở 99,8\% lượt gọi của cấu hình mặc định; các lượt còn lại quay về thứ tự RRF như thiết kế.

\subsection{Ảnh hưởng của tìm kiếm web}
\label{sec:ablation_web}

Ở chế độ tự động, Agent gần như không gọi web (khoảng 1 lần trên 100 câu) vì kho nội bộ đã bao phủ bộ câu hỏi, nên kết quả rất gần chế độ tắt. Chế độ luôn tìm không làm tăng độ đúng nhưng làm p50 tăng gấp đôi (khoảng 6 giây lên 12--15 giây), số token mỗi câu tăng khoảng 56\% (khoảng 22.100 lên 34.400--34.900), chi phí mỗi lượt đo tăng từ khoảng 0,24 USD lên 0,34--0,35 USD, và số câu phải kết thúc bằng lượt \code{finish} bắt buộc tăng lên 18--22 câu (so với 0--1 câu ở các chế độ khác) do Agent tốn thêm lượt cho web. Độ bám nguồn ở chế độ này vẫn cao (98,4--99,5\%). Như vậy câu trả lời cho RQ3 là để Agent tự quyết định khi nào tìm web là lựa chọn tốt hơn; chế độ luôn tìm chỉ nên dùng khi người chơi cần thông tin mới nhất, ví dụ thay đổi của một bản vá.

\subsection{Ảnh hưởng của kiểu truy xuất}
\label{sec:ablation_retrieval}

Ở trọn luồng, khác biệt giữa dense, BM25 và hybrid nhỏ hơn nhiều so với khi đo truy xuất riêng, vì Agent có thể tìm lại nhiều lần với truy vấn khác nhau và tự kiểm tên, bù được phần nào điểm yếu của từng phương pháp. Dù vậy, hybrid kết hợp L12 vẫn cho độ đúng cao nhất, đặc biệt ở câu khó (97,8\%).

% ==============================================================================
\section{Hiệu năng và chi phí}
\label{sec:perf}

\subsection{Độ trễ}
\label{sec:latency}

Với cấu hình mặc định, độ trễ tới câu trả lời đầu tiên có p50 6,6 giây và p95 15,1 giây. Trong một lượt, phần lớn thời gian thuộc về các lượt gọi LLM: vòng lặp Agent cần vài lượt gọi mô hình (tìm kiếm, mở bằng chứng, kết thúc), sau đó là bước sinh câu trả lời. Truy xuất lai chỉ mất khoảng 67 ms mỗi truy vấn, bộ xếp hạng lại cục bộ khoảng 0,9 giây. Do câu trả lời chỉ được gửi sau khi kiểm tra xong, thời gian tới token đầu tiên trên web gần bằng thời gian tới câu trả lời đầy đủ; giao diện bù lại bằng các sự kiện trạng thái trong lúc xử lý.

Nhóm đã xây dựng kịch bản kiểm thử tải bằng Locust (đo thời gian tới token đầu tiên và tổng thời gian của luồng SSE với 20 và 50 người dùng đồng thời) và các script tính p50/p95 theo từng bước từ trace, nhưng tại thời điểm viết báo cáo chưa chạy kiểm thử tải với chế độ RAG thật trên máy chủ triển khai, vì mỗi lượt tốn chi phí API đáng kể. Vì vậy báo cáo chưa có số đo cho yêu cầu chịu tải 50 người dùng đồng thời.

\subsection{Chi phí}
\label{sec:cost}

Chi phí DeepSeek được đo bằng chênh lệch số dư tài khoản trước và sau mỗi lượt 200 câu. Với cấu hình mặc định, một lượt tốn khoảng 0,24 USD, tức khoảng 0,0012 USD mỗi câu. Nếu ngoại suy cho mức 1.000 câu mỗi ngày, chi phí LLM khoảng 1,2 USD/ngày hay khoảng 36 USD/tháng, cao hơn một chút so với mục tiêu 30 USD/tháng. Đây là ước tính thô vì bộ câu đánh giá có tỉ lệ câu khó cao hơn câu hỏi thực tế, và chưa tính chi phí Tavily (rất ít khi dùng ở chế độ tự động). Để kiểm soát chi phí trong vận hành, hệ thống áp dụng hạn mức câu hỏi và token theo người dùng cùng trần token chung mỗi ngày (Mục~\ref{sec:impl_api}).

% ==============================================================================
\section{Triển khai hệ thống}
\label{sec:deploy}

\subsection{Môi trường triển khai}
\label{sec:deploy_env}

Hệ thống được triển khai trên một VPS 4 vCPU, 8 GB RAM, chạy Ubuntu 24.04 và Docker Compose, tại tên miền của nhóm \url{https://tarnished.liemdai.io.vn} (tại thời điểm viết báo cáo). Cấu hình production là một tệp overlay chồng lên tệp Compose gốc, đặt cố định chế độ RAG, lưu hội thoại vào PostgreSQL, không công khai cổng cơ sở dữ liệu, thêm kiểm tra sức khỏe và giới hạn tài nguyên cho từng container (Bảng~\ref{tab:resources}). Ngoài môi trường production còn có môi trường staging cùng cấu hình để thử trước khi phát hành.

\begin{table}[H]
    \centering
    \caption{Giới hạn tài nguyên các container ở môi trường production}
    \label{tab:resources}
    \small
    \begin{tabular}{|l|r|r|l|}
        \toprule
        \textbf{Container} & \textbf{RAM} & \textbf{CPU} & \textbf{Vai trò} \\
        \midrule
        api & 3 GB & 2,5 & FastAPI, Agent, E5 và MiniLM-L12 ONNX \\ \hline
        qdrant & 1 GB & 1,5 & Chỉ mục corpus và ký ức \\ \hline
        postgres & 1 GB & 1,0 & Dữ liệu quan hệ \\ \hline
        redis & 192 MB & 0,5 & Giới hạn tần suất, hạn mức \\ \hline
        bot & 192 MB & 0,5 & Discord bot \\ \hline
        caddy & 256 MB & 0,5 & HTTPS, web tĩnh, reverse proxy \\
        \bottomrule
    \end{tabular}
\end{table}

\subsection{Bảo mật vận hành}
\label{sec:security}

\begin{itemize}
    \item Mạng: chỉ mở cổng 22 (SSH, chỉ đăng nhập bằng khóa), 80 và 443. Caddy tự cấp và gia hạn chứng chỉ HTTPS, thêm các header bảo mật (HSTS, CSP), chặn các đường dẫn tài liệu API và endpoint dành cho phát triển ở môi trường production, giới hạn kích thước thân yêu cầu 64 KB.
    \item Container: api, bot và Caddy chạy với \code{no-new-privileges} và bỏ toàn bộ capability (Caddy chỉ thêm lại quyền mở cổng thấp).
    \item Xác thực: bot dùng token dịch vụ; trang quản trị dùng token quản trị hoặc phiên của tài khoản Discord nằm trong danh sách quản trị; phiên người dùng là cookie \code{httpOnly} lưu dạng băm trong PostgreSQL; ứng dụng từ chối khởi động ở môi trường production nếu còn giá trị mặc định \code{change-me}.
    \item Chống lạm dụng: giới hạn tần suất theo phút và hạn mức theo ngày cho mọi endpoint hỏi đáp, Redis lỗi thì từ chối thay vì cho qua.
    \item Riêng tư: API và Caddy không ghi access log, nên địa chỉ IP người dùng không bị lưu; log ứng dụng không ghi nội dung hội thoại.
\end{itemize}

\subsection{Tích hợp liên tục và quy trình phát hành}
\label{sec:ci}

Kho mã dùng quy trình nhánh: mỗi tính năng làm trên nhánh riêng tách từ \code{develop}, mở pull request về \code{develop} và được review trước khi hợp nhất; \code{main} chứa bản phát hành. Workflow GitHub Actions chạy cho mỗi pull request với ba nhóm việc: backend (ruff kiểm tra và định dạng, kiểm tài liệu OpenAPI khớp mã, unit test, đo độ bao phủ module RAG lõi với ngưỡng 60\%, integration test với PostgreSQL, Qdrant, Redis thật), web (kiểm định dạng và build) và bot (kiểm định dạng và kiểm kiểu TypeScript). Việc triển khai dùng script \code{deploy.sh}: kiểm tra cây làm việc sạch và tệp \code{.env} hợp lệ, sao lưu PostgreSQL trước khi triển khai, kéo mã, build lại container và chờ API báo khỏe; rollback bằng cách chạy lại script với tag hoặc commit cũ. PostgreSQL được sao lưu hằng ngày (giữ 7 bản) và Qdrant được chụp snapshot hằng tuần (giữ 4 bản), có script khôi phục kiểm tra bản sao lưu trước khi thay cơ sở dữ liệu đang chạy.

\subsection{Kịch bản demo}
\label{sec:demo}

Nhóm chuẩn bị kịch bản demo năm câu, mỗi câu thể hiện một năng lực của Agent:
\begin{enumerate}
    \item Câu lore tiếng Việt: ``Ranni là ai?'' -- trả lời có trích dẫn, xem đoạn nguồn bằng popover.
    \item Câu gõ sai, không dấu: ``chia stoneswrod dung voi tuong nao'' -- Agent kiểm tên rồi tìm đúng trang Stonesword Key.
    \item Câu nhiệm vụ nhiều bước -- câu trả lời dạng danh sách đánh số, trích dẫn ở từng bước.
    \item Câu build có tính toán: ``Build Samurai dùng Uchigatana ở level 60'' -- bảng tám chỉ số tính bằng \code{run\_python}, nêu rõ giả định.
    \item Câu cần hỏi lại -- Agent hỏi thông tin còn thiếu, người chơi trả lời, hệ thống tiếp tục.
\end{enumerate}
Cùng các câu được thử trên Discord bằng \code{/ask} và bằng cách nhắc tên bot, và trang quản trị được dùng để mở trace của từng câu. Ảnh chụp các bước demo trên web và Discord được trình bày ở Mục~\ref{sec:impl_web_ui} và Mục~\ref{sec:impl_bot}.

\subsection{Hướng dẫn cài đặt và sử dụng}
\label{sec:guide}

Hướng dẫn cài đặt môi trường phát triển và tái lập đánh giá được trình bày ở Phụ lục~\ref{app:install}. Với người chơi, cách sử dụng rất đơn giản: trên web, nhập câu hỏi bằng tiếng Việt hoặc tiếng Anh vào ô chat, bấm vào số trích dẫn để xem đoạn nguồn và mở trang wiki gốc, chọn chế độ web OFF/ON/ALWAYS nếu cần; đăng nhập bằng Discord để giữ lịch sử trên nhiều thiết bị. Trên Discord, mời bot vào server, dùng \code{/ask <câu hỏi>} hoặc nhắc tên bot; quản trị viên server có thể dùng \code{/setup} để bot trả lời mọi tin nhắn trong một kênh.

% ==============================================================================
\section{Thảo luận}
\label{sec:discussion}

Điểm mạnh. Hệ thống đạt độ đúng 97,8\% và độ bám nguồn 96,2\% trên bộ 200 câu, vượt xa yêu cầu đặt ra; xử lý tốt câu tiếng Việt và câu gõ sai trên kho tri thức tiếng Anh; mọi câu trả lời có dữ kiện đều kèm trích dẫn được kiểm nguyên văn; không câu nào vượt hạn 45 giây. Việc tách rõ Agent (quyết định ngữ nghĩa) và máy chủ (kiểm tra tất định) giúp hệ thống vừa linh hoạt vừa an toàn: Agent có thể tìm lại, kiểm tên, tính toán, còn máy chủ đảm bảo không trích dẫn sai và không bịa điều kiện người chơi.

Hạn chế.
\begin{itemize}
    \item Hành vi với câu mơ hồ còn yếu (behavior pass 70,6\%): Agent thường chọn một cách hiểu thay vì hỏi lại.
    \item Độ trễ p95 (15,1 giây) chưa đạt mục tiêu 10 giây; token đầu vào của Agent lớn (khoảng 22.000 token mỗi câu) do mỗi lượt gửi lại toàn bộ lịch sử công cụ.
    \item Kiểm tra trích dẫn chỉ ở mức hình thức; máy chủ không kiểm lại phép tính build do Agent thực hiện mà dựa vào các câu \code{assert} Agent tự viết.
    \item Kho tri thức còn thiếu một số bảng số liệu có cấu trúc (chỉ số vũ khí theo cấp nâng), giới hạn khả năng trả lời câu so sánh chi tiết.
    \item Trên Discord, bảng Markdown trong câu trả lời (ví dụ bảng chỉ số build) hiển thị ở dạng ký tự vì embed không hỗ trợ bảng.
    \item Hệ thống phụ thuộc API bên ngoài (DeepSeek, Tavily); chưa có kiểm thử tải thật.
\end{itemize}

Các yếu tố ảnh hưởng tới tính hợp lệ của đánh giá.
\begin{itemize}
    \item Kích thước tập đo: 183 câu dữ kiện nên mỗi câu tương ứng 0,55 điểm phần trăm; khác biệt nhỏ giữa các cấu hình chưa có ý nghĩa thống kê, và mỗi cấu hình mới chạy một lần.
    \item Người chấm là LLM: dù đã đối chiếu với một mô hình chấm khác (khớp 181/183 câu) và dùng mô hình khác với mô hình được đánh giá, việc chấm bằng LLM vẫn có thể có thiên lệch~\cite{zheng2023judge}; chưa có lượt chấm tay quy mô lớn.
    \item Bộ câu do nhóm tự xây dựng: câu hỏi và nguồn gold do nhóm viết dựa trên chính kho tri thức, nên có thể dễ hơn câu hỏi thực tế của người chơi; lượt rà soát cuối có sự hỗ trợ của trợ lý AI.
    \item Điều kiện đo độ trễ: đo tuần tự trên máy phát triển, không có tải đồng thời; độ trễ còn phụ thuộc tình trạng API DeepSeek tại thời điểm đo.
\end{itemize}

% ==============================================================================
\section{Tóm tắt chương}
\label{sec:ch4_summary}

Chương này đã đánh giá hệ thống trên bộ 200 câu song ngữ với 11 cấu hình. Truy xuất lai kết hợp bộ xếp hạng lại MiniLM-L12 cục bộ đạt Recall@5 86,6\% và MRR@5 0,739. Cấu hình mặc định đạt độ đúng 97,8\%, faithfulness 96,2\%, p50 6,6 giây, không câu nào vượt 45 giây; điểm yếu chính là hành vi với câu mơ hồ. Thí nghiệm ablation cho thấy bộ xếp hạng lại có ích nhất ở câu khó, còn chế độ luôn tìm web làm tăng gấp đôi độ trễ mà không tăng độ đúng. Hệ thống đã được triển khai trên VPS bằng Docker Compose với HTTPS, CI và các lớp bảo vệ vận hành.
